% Body-only replacement for Section 3.3.3. Keep the existing subsubsection title/label.
% Uses the project natbib citation style. Do not add a second bibliography here.
% Four paragraph headings; eight unique works.

Sustained interaction requires models to preserve both speakers' timing and update responses as new speech arrives. Native dynamic interaction makes speech, silence, or listening--speaking transitions part of the content model's predictions; decoding and playback modules execute the output. Multi-stream modeling preserves parallel channels, flattening serializes their contents, and chunk-based alternation organizes temporal progression. These choices can be combined rather than treated as mutually exclusive architectures~\citep{DBLP:journals/corr/abs-2410-00037,DBLP:conf/acl/ZhangCDCWZLYTDZ25,DBLP:conf/emnlp/VeluriPYGG24}.

\paragraph{Multi-stream parallel modeling.}
Separate speech channels retain pauses and overlaps without imposing turn boundaries. The early dGSLM uses cross-attending Transformer towers to predict speech units and durations, with evaluation based on offline dialogue continuation~\citep{DBLP:journals/tacl/NguyenKCAHETASM23}. Moshi aligns user speech, assistant text, and assistant speech on an 80~ms grid. Training jointly models both speakers; online inference observes the real user channel and generates assistant-side content. Near-silent audio and text padding keep the sequence advancing during silence. A Temporal Transformer models cross-frame history, while a smaller Depth Transformer predicts within-frame codebooks; locally preceding text provides the Inner Monologue condition~\citep{DBLP:journals/corr/abs-2410-00037}. Fun-Audio-Chat-Duplex further separates a 5~Hz backbone from a 25~Hz speech refinement head. Both designs offload local acoustic prediction, while the latter also reduces backbone update frequency~\citep{DBLP:journals/corr/abs-2512-20156}.

\paragraph{Flattened sequence modeling.}
OmniFlatten reuses a standard GPT-style backbone by repeatedly serializing user speech, assistant text, and assistant speech chunks. Chunk-level loose alignment and silence tokens retain inactive intervals; real user chunks are inserted between assistant generations. Unlike Moshi's hierarchical multi-stream prediction, all channel tokens occupy positions in one backbone sequence. Ablations favor the progression from modality alignment through half-duplex to full-duplex learning, and removing assistant text reduces response quality. The model is evaluated on turn-taking and interruption-induced stopping, but the paper explicitly excludes user and assistant backchannels~\citep{DBLP:conf/acl/ZhangCDCWZLYTDZ25}.

\paragraph{Chunk-based synchronization.}
SyncLLM alternates fixed-duration speech chunks, deduplicates speech units, and retains boundaries through periodic speaker markers. While an assistant chunk plays, the concurrent user chunk is incomplete. The model estimates that user chunk to prepare the next assistant chunk, then replaces the estimate in its history when observed input arrives. This trades input waiting for prediction error; its synchronization evidence comes from two-model dialogue simulation~\citep{DBLP:conf/emnlp/VeluriPYGG24}. Internal alternation still permits overlapping speech in real time. Output synchronization poses a separate constraint: MiniCPM-o 4.5 uses Time-Aligned Interleaving (TAIL) to adjust text generation against cumulative playback progress, limiting how far text runs ahead of speech. SyncLLM addresses incomplete incoming speech, whereas TAIL coordinates outgoing text and audio~\citep{DBLP:journals/corr/abs-2604-27393}.

\paragraph{State and timing learning.}
Explicit state tokens make transitions trainable. SALMONN-\allowbreak omni conditions on environmental and assistant-history audio, using \texttt{<shift>} for listening--speaking transitions and \texttt{<think>} for listening or waiting intervals. Preference optimization reduces unnecessary stopping for user acknowledgments~\citep{DBLP:conf/nips/YuWYCTZSLWZ25}. Within a flattened sequence, BayLing-Duplex distinguishes silence, response onset, text completion, and completion of the full output. Observed user chunks condition the next output window with the same causal shift during training and inference; its 0.8~s blocks consequently constrain when an interruption can alter generation. Reweighting silence and response-boundary tokens improves turn-taking, while timing-based preference pairs keep response text fixed and vary onset or stopping times. Its interruption tests assess both stopping and the answer to the new question~\citep{DBLP:journals/corr/abs-2606-14528}. Temporal representations therefore require behavior-specific supervision.

Native approaches let content and timing share an evolving dialogue history. Finer updates create more opportunities to react but require more frequent encoding and generation; larger chunks delay intervention. Audible response also depends on decoding and playback buffers. Comparisons should therefore separate turn-taking, stopping, and content revision, and identify whether evidence comes from prerecorded input, model simulation, or human interaction~\citep{DBLP:journals/corr/abs-2410-00037,DBLP:conf/emnlp/VeluriPYGG24,DBLP:journals/corr/abs-2606-14528}.
